% !TEX root = ../notes.tex

\documentclass[../notes.tex]{subfiles}

\begin{document}

\section{May 6}
Here we go.

\subsection{The Cohomology of the Upper-Half Plane}
We now use \Cref{prob:bc-as-quotient}.
\begin{definition}[Drinfeld upper-half plane]
	Fix a finite field extension $F$ of $\QQ_p$, and let $C$ be an algebraically closed perfectoid field containing $F$. Then we define $\mc H^{1,\diamond}_{F,C}$ to be $\PP^1(C)\setminus\PP^1(F)$.
\end{definition}
\begin{remark}
	There is a natural action of $\op{GL}_2(F)$ on $\PP^1(C)$, which descends to $\mc H^{1,\diamond}_{F,C}$.
\end{remark}
\begin{remark}
	By \Cref{prob:bc-as-quotient}, we see that
	\[\op{BC}(-1/t)^*\times_{\op{Spd}k}\op{Spd}C=\mc H^{1,\diamond}_{F,C}/F.\]
	Here, $F$ acts on $\mc H^{1,\diamond}_{F,C}$ as the restriction of the action of $\op{GL}_2(F)$ to the matrices $\begin{bsmallmatrix}
		1 & x \\ & 1
	\end{bsmallmatrix}$.
\end{remark}
Thus, we have a fairly explicit presentation of the punctured Banach--Colmez space, and we see that we are tasked with computing
\[\mathrm R\Gamma(\mc H^{1,\diamond}_{F,C}/F;\OO^\flat).\]
We will do this calculation in stages.
\begin{lemma} \label{lem:upper-half-plane-cohomology}
	Fix a finite field extension $F$ of $\QQ_p$. Then
	\[\mathrm H^i_c\left(\mc H^{1,\diamond}_{F,C};\OO^\flat\right)=\begin{cases}
		C^\infty\left(\PP^1(F);C^\flat\right)/C^\flat & \text{if }i=1, \\
		C^\flat & \text{if }i=2, \\
		0 & \text{otherwise}.
	\end{cases}\]
	Here, $C^\infty$ means locally constant.
\end{lemma}
\begin{proof}
	Note that $\PP^1(C)=\mc H^{1,\diamond}_{F,C}\sqcup\PP^1(F)$ is a disjoint union into an open subset and its closed complement. Thus, some formal nonsense with six-functor formalisms provides a fiber sequence
	\[\mathrm R\Gamma_c\left(\mc H^{1,\diamond}_{F,C};\OO^\flat\right)\to\mathrm R\Gamma_c\left(\PP^1(C);\OO^\flat\right)\to\colim_i\mathrm R\Gamma(U_i;\OO^\flat),\]
	where $\{U_i\}_i$ are some open disks which have intersection $\PP^1(F)$. The right-hand complex is basically a formal completion of $\PP^1(F)$; note that we are basically doing some of \v{C}ech cohomology. It turns out that the colimit is an Eilenberg--Maclane space of the group $C^\infty(\PP^1(F),C^\flat)$. As such, we have a long exact sequence in homotopy groups
	\[0\to\mathrm H^0_c\left(\mc H^{1,\diamond}_{F,C};\OO^\flat\right)\to\mathrm H^0\left(\PP^1_C;C^\flat\right)\to C^\infty\left(\PP^1(F);C^\flat\right)\to\mathrm H^1_c\left(\mc H^{1,\diamond}_{F,C};\OO^\flat\right)\to\underbrace{\mathrm H^1_c\left(\PP^1(C);\OO^\flat\right)}_0\to0.\]
	Thus, we see that $\mathrm H^0_c\left(\mc H^{1,\diamond}_{F,C};\OO^\flat\right)=0$, and the statement for degree one also follows. A similar exact sequence allows us to compute
	\[\mathrm H^2_c\left(\mc H^{1,\diamond}_{F,C};\OO^\flat\right)\cong\mathrm H^2_c\left(\PP^1(C);\OO^\flat\right)=C^\flat.\]
	After this point, the cohomology of $\PP^1(C)$ vanishes, so we see that all higher cohomology of $\mc H^{1,\diamond}_{F,C}$ vanishes. 
\end{proof}
\begin{remark}
	Our group $\op{GL}_2(F)$ acts naturally on $C^\infty\left(\PP^1(F);C^\flat\right)/C^\flat$ and trivially on $C^\flat$.
\end{remark}
\begin{remark}
	The representation $C^\infty\left(\PP^1(F);C^\flat\right)$ of $\op{GL}_2(F)$ is known as the Steinberg representation.
\end{remark}
Our next task is to take the quotient by $F$. It will be convenient to remember the torus, so we let $B(F)\subseteq\op{GL}_2(F)$ denote the subgroup of diagonal matrices, and let $T(F)$ be the subgroup of diagonal matrices. To this end, note that we have a sequence
\[\mc H^{1,\diamond}_{F,C}/B(F)\stackrel g\to\mathrm{pt}/B(F)\stackrel h\to\mathrm{pt}/T(F).\]
We would like to calculate $(hg)_*\OO^\flat$, which is $\mathrm R\Gamma\left(\mc H^{1,\diamond}_{F,C}/F\right)$ with a residual action by $T(F)$.

To this end, we want to know something about the six functor formalism. For a morphism $f\colon X\to Y$, the functor $f^*$ admits a left adjoint $f_*$. There is also a dualizing object $f^!C^\flat$, and if $f$ is smooth, $f^!(-)=f^*(-)\otimes f^!C^\flat$. The functor $f^!$ also admits a left adjoint $g_!$. Furthermore, if $f$ is smooth, $f^*$ has left adjoint $f_\natural\coloneqq f_!\left(-\otimes f^!C^\flat\right)$. Applied this formalism to our situation, we know $g$ and $h$ are smooth, so
\[(h\circ g)_*\OO^\flat=\left((hg)_\natural\OO^\flat\right)^\lor,\]
where the dual has appeared because we are undoing two adjoints. Let's start with $g_\natural$.
\begin{lemma}
	Fix everything as above. Then the complex $g_\flat\OO^\flat$ has cohomology $C^\infty\left(\PP^1(F),C^\flat\right)/C^\flat$ in degree $-1$, $C^\flat$ in degree zero, and it vanishes everywhere else.
\end{lemma}
\begin{proof}
	This follows by applying Poincar\'e duality to \Cref{lem:upper-half-plane-cohomology}. Note that we get a shift because the dualizing object (pulled back from $\PP^1$) is trivial but shifted by $-2$.
\end{proof}
\begin{lemma}
	Fix everything as above. Then $h_\natural C^\flat=C^\flat$.
\end{lemma}
\begin{proof}
	By passing to homology, the main point is to show that
	\[\mathrm H_i(F;C^\flat)=\begin{cases}
		C^\flat & \text{if }i=0, \\
		0 & \text{otherwise}.
	\end{cases}\]
	To this end, we filter the space $F$ as $\bigcup_{m\ge0}\varpi^{-m}\OO_F$, but of course, all the pieces $\varpi^{-m}\OO_F$ are abstractly isomorphic to $\ZZ_p^d$, and the transition maps are given by multiplication by $\varpi$ (to the power of the degree). Now, $\mathrm H_i(\ZZ_p^d;C^\flat)$ is the homology of $\mathrm B\ZZ_p^{d}$ with coefficients in $C^\flat$. But this is the same as the cohomology of $\left(S^1_p\right)^{d}$ with coefficients in $C^\flat$, so we will find homology only in degrees $0$ and $1$.
	
	From here, the homology in degree zero survives the transition maps, but the homology in degree one does not survive: we find that we are calculating a colimit of the diagram
	\[C^\flat\stackrel\varpi\to C^\flat\stackrel\varpi\to\cdots,\]
	which vanishes.
\end{proof}
\begin{lemma}
	Fix everything as above. Then
	\[h_\natural\left(C^\infty\left(\PP^1(F),C^\flat\right)/C^\flat\right)=\Sigma^d\chi,\]
	where $d\coloneqq[F:\QQ_p]$ and $\chi\colon T(F)\to C^\flat$ is the character $\op{diag}(t_1,t_2)\mapsto\varepsilon\op N_{F/\QQ_p}(t_1/t_2)$, where $\varepsilon\colon\QQ_p^\times\to\FF_p^\times$ sends $p\mapsto1$ and projects $\ZZ_p^\times\onto\FF_p^\times$.
\end{lemma}
\begin{proof}
	Let $s$ be the quotient $\mathrm{pt}/T(F)\to\mathrm{pt}/B(F)$. Then $s_!C^\flat$ consists of the locally constant functions $B(F)/T(F)\to C^\flat$, which means that we are looking at locally constant functions $F\to C^\flat$, which is the same as $C^\infty\left(\PP^1(F),C^\flat\right)/C^\flat$, where the idea is that $C^\flat$ can be used to normalize the output of $\infty\in\PP^1(F)$ to zero.

	Thus,
	\[h_!\left(C^\infty\left(\PP^1(F),C^\flat\right)/C^\flat\right)=h_!s_!C^\flat,\]
	but this is just $C^\flat$ because $hs$ is the identity. We now want to compute
	\[h_\natural\left(C^\infty\left(\PP^1(F),C^\flat\right)/C^\flat\right)=h_!\left(C^\infty\left(\PP^1(F),C^\flat\right)/C^\flat\otimes h^!C^\flat\right),\]
	so the main point is to compute $h^!C^\flat$. Well, for any $p$-adic Lie group $G$ over $\QQ_p$, the dualizing sheaf is $\Sigma^{\dim_{\QQ_p}G}\det\mf g$, where $\mf g$ is the adjoint representation. Comparing the two modular characters of $B(F)$ and $T(F)$ (and their dimensions) produces the result.
\end{proof}
In total, once combined with \Cref{prop:drinfeld}, we have calculated that $\mathrm R\Gamma\left(\mc H^{1,\diamond}_{F,C}/F\right)^{h\mc C_q}$ has homotopy groups $k$ in degree zero and $k$ in degree $d+1$ (and it vanishes everywhere).\footnote{Notably, taking $h_\natural$ will not necessarily split the complex $g_\natural\OO^\flat$. Nonetheless, we can still use this to compute homotopy groups because we have found that $g_\natural\OO^\flat$ lives in a fiber bundle of Eilenberg--Maclane spaces, so the same will be true after $h_\natural$.}
\begin{remark}
	We may be interested in some residual action: namely, recall that $\mathrm R\Gamma\left(\op{Spd}\widehat L/G^\circ_h\right)$ should have a residual action by $\ZZ_p^\times$ because $G^\circ_h$ is the kernel of the determinant map $G_h\to\ZZ_p^\times$. By tracing through all the isomorphisms, we see that $\ZZ_p^\times$ acts trivially in degree zero, and $\alpha\in\ZZ_p^\times$ acts by $\overline\alpha^{h-1}$ (in $\FF_p^\times$) in degree $h$. This is nontrivial as soon as $p>2$.
\end{remark}

% A natural place to start would be to compute $\mathrm R\Gamma(\mc H^{1,\diamond}_{F,C};\OO^\flat)$ while remembering the $F$-action. To this end, we recall that $\mc H^{1,\diamond}_{F,C}$ embeds as an open subset of $\PP^1_C$, and the complement is the proper subset $\PP^1(F)$. Thus, there is a version of excision in our six functor formalism, which provides a long exact sequence
% \[0\to\mathrm H^0_c(\mc H^{1,\diamond}_{F,C})\to\mathrm H^0(\PP^1_C)\to\op{LC}\left(\PP^1(F),C\right)\to\mathrm H^1_c(\mc H^{1,\diamond}_{F,C})\to\cdots.\]
% This allows us to compute compactly supported cohomology, and then we pass to cohomology via Poincar\'e duality. We will discuss this next time.

\subsection{Application to Chromatic Splitting}
We are now ready to return to homotopy theory. Let's repeat the argument of \Cref{rem:get-chromatic-splitting}.
\begin{theorem}
	The chromatic splitting conjecture is true for primes $p$ and heights $h$ with $(p-1)\nmid(h-1)$.
\end{theorem}
\begin{proof}
	There is a composite
	\[E_{h-1}\to E(L)\to E(\widehat L),\]
	where we recall that $E_{h-1}=E(k;\Gamma_{h-1})$. Here, $E_{h-1}$ has a remaining action by $G_{h-1}$, but the other two objects have natural actions by both $G_{h-1}$ and $G_h$. We can further give $E_{h-1}$ a trivial action by $G_{h-1}$, which makes the composite fully equivariant. For example, we have received a map
	\[E_{h-1}\to E(\widehat L)^{hG^\circ_h},\]
	which is $G_{h-1}$-equivariant. We can take further invariants to see that there is a map
	\[E_{h-1}\to\left(E(\widehat L)^{hG^\circ_h}\right)^{h\FF_p^\times}.\]
	Taking this$\pmod{p,v_1,\ldots,v_{h-2}}$, we get  $K(n)$ on the left. On the right, we have some spectral sequence starting with $\FF_p^\times$-fixed points of our previous cohomology calculation. Thus, if $p-1\nmid h-1$, we find that $\FF_p^\times$ kills the part in different degree, so we find that the displayed map is a $G_{h-1}$-equivariant isomorphism. This is what we needed for chromatic splitting!
\end{proof}
\begin{remark}
	Even when $p-1\mid h-1$, we can measure how badly we failed. For example, we may let $C$ be the cofiber of the natural map
	\[E_{h-1}\to\left(E(\widehat L)^{hG^\circ_h}\right)^{h\FF_p^\times}.\]
	By our homotopy group calculation, arguing as in \Cref{rem:get-chromatic-splitting}, we see that $C$ is $\Sigma^hE_{h-1}$ to recover the other homotopy group. One could now try to understand $C$ more finely. For example, by Galois theory, $K(h-1)$-local spectra have the same data as a $G_{h-1}$-equivariant $E_{h-1}$-module, so we can view $C$ as an invertible element in the category of $K(h-1)$-local spectra.
\end{remark}

\end{document}